% =====================================================================================================
% IEEE conference paper draft (v2). The previous draft is kept in main_v1_old.tex.
%
% NO RESULT IS TYPED INTO THIS FILE. Every number is \R{key}, read from generated/numbers.tex; every table
% and figure is read from generated/. Both are produced by notebooks/er_full_pipeline_colab.ipynb and
% copied here with:   python scripts/import_results.py <er_results_RUN.zip>
%
% Until results are imported the PDF compiles with red placeholders: [CE.ap] marks a missing number and
% [claim B pending] marks wording that depends on the paper-readiness check.
%
% \IfClaim{X}{supported wording}{unsupported wording} picks the sentence from the readiness flag:
%   B = validation-selected entity operating point beats the pairwise one
%   C = hard veto lowers AP of the strongest matcher
%   D = veto precision falls as matcher quality rises
%   E = conflict evidence helps as a soft feature
%   F = conflict-aware review routing beats uncertainty routing
%   P = cross-encoder reproduces the published RoBERTa results
%   R = the DeBERTa-v3 cross-encoder agrees with RoBERTa on claims C and F (replication)
% After importing, read every selected branch: the wording is chosen automatically, but how strongly a
% result is phrased remains the authors' decision.
%
% Before submission: complete the Acknowledgment. IEEE requires any use of AI-generated content
% (text, code, figures) to be disclosed there; describe honestly what was AI-assisted.
%
% Page budget: IEEE BigData regular papers allow 10 pages, many IEEE workshops 6-8. To reach 6 pages,
% move table_robustness, table_cost and the reliability figure to an appendix or the repository.
% =====================================================================================================
\documentclass[conference]{IEEEtran}

\usepackage{cite}
\usepackage{amsmath,amssymb,amsfonts}
\usepackage{graphicx}
\usepackage{booktabs}
\usepackage{xcolor}
\usepackage{url}
\usepackage{textcomp}

% ---- result plumbing ---------------------------------------------------------------------------------
\newcommand{\R}[1]{\ifcsname res:#1\endcsname\csname res:#1\endcsname\else{\color{red}\textbf{[#1]}}\fi}
\newcommand{\IfClaim}[3]{\ifcsname res:flag.#1\endcsname
  \ifnum\csname res:flag.#1\endcsname=1 #2\else #3\fi
  \else{\color{red}[claim #1 pending]}~#2\fi}
\newcommand{\GenTable}[2]{\IfFileExists{generated/#1.tex}{\input{generated/#1.tex}}{%
  \begin{table}[t]\centering\fbox{\parbox{0.92\columnwidth}{\centering\color{red}\textbf{#2}\\
  pending: \texttt{\detokenize{generated/#1.tex}}}}\end{table}}}
\newcommand{\GenFig}[5]{% env, file, width, caption, label
  \begin{#1}[t]\centering
  \IfFileExists{generated/#2.pdf}{\includegraphics[width=#3]{generated/#2.pdf}}{%
    \fbox{\parbox{0.9\columnwidth}{\centering\color{red} pending: \texttt{\detokenize{generated/#2.pdf}}}}}
  \caption{#4}\label{#5}\end{#1}}
\IfFileExists{generated/numbers.tex}{\input{generated/numbers.tex}}{}

\begin{document}

\title{\IfClaim{B}{Tune for Entities, Not Pairs: Operating Points, Attribute Constraints and Review
Routing for Product Entity Resolution}{Attribute Constraints, Matcher Quality and Review Routing in
Product Entity Resolution: An Empirical Study}}

\author{\IEEEauthorblockN{[Author 1], [Author 2], [Author 3]}
\IEEEauthorblockA{\textit{Department of Computer Science and Engineering}\\
\textit{SRM Institute of Science and Technology}\\
[City], India\\
\{[email1], [email2]\}@srmist.edu.in}
\and
\IEEEauthorblockN{[Guide name]}
\IEEEauthorblockA{\textit{Department of Computer Science and Engineering}\\
\textit{SRM Institute of Science and Technology}\\
[City], India\\
[email]}}

\maketitle

% ======================================================================================================
\begin{abstract}
Most evaluations of product matching stop at pairwise F1, measured at whichever threshold maximises it.
Catalogue owners, however, work with the merged entities, and one wrong link can join two product
families. This paper examines three choices made after a matcher has scored its candidate pairs: where
to put the threshold, whether to overrule the matcher when two offers state conflicting attributes, and
which pairs to hand to a human reviewer when only a few can be checked. We run eight matchers from four
architecture families (string and feature-based baselines, a bi-encoder and two cross-encoders based on
RoBERTa and DeBERTa-v3, all fine-tuned, and zero-shot LLMs from two model families) on all nine variants of WDC
Products and on \R{nextra} Magellan benchmarks. Thresholds, calibrators and other settings are fitted
on validation data only, and every claim is judged against success criteria that we fixed before running
the experiments.
\IfClaim{B}{Picking the threshold that maximises entity-level F1 on validation lifts the cross-encoder's
entity F1 from \R{CE.entpw} to \R{CE.entbest}, with a gain in \R{CE.entbestwins} variants.}{Picking the
threshold that maximises entity-level F1 on validation gives no consistent gain over the pairwise
threshold (better in \R{CE.entbestwins} variants for the cross-encoder).}
\IfClaim{D}{A hard attribute veto becomes less precise as the matcher gets stronger (Spearman
$\rho=\R{rho}$ over \R{npts} runs). With the cross-encoder, only \R{CE.vetoprec} of the accepted pairs it
rejects are non-matches, so in most cases the veto undoes a correct merge.}{How precise a hard attribute
veto is does not depend consistently on matcher quality ($\rho=\R{rho}$).}
\IfClaim{F}{When the conflict signal is used to pick pairs for review instead of rejecting them, it adds
\R{CE.awareunc} F1 over uncertainty-based selection at a \R{budget} budget.}{Using the conflict signal to
pick pairs for review brings no gain over selecting them by calibrated uncertainty.}
As a check on our protocol, the cross-encoder stays within \R{CE.paritybias} F1 of the published RoBERTa
results on average.
\end{abstract}

\begin{IEEEkeywords}
entity resolution, product matching, operating point, calibration, attribute constraints, human review
\end{IEEEkeywords}

% ======================================================================================================
\section{Introduction}

Online marketplaces receive the same product from many sellers, and each seller writes its own title.
One may list \texttt{Apple iPhone 15, 128GB, Black}, another \texttt{iPhone 15 Apple 128 GB Midnight};
both describe one phone. Compare \texttt{Apple iPhone 15, 128GB} with \texttt{Apple iPhone 15 Plus
128GB}: these two share even more words, yet they are different models. Entity resolution (ER) has to join the
first pair and keep the second apart.

Research on entity matching usually reports precision, recall and F1 over the candidate pairs released
with a benchmark~\cite{peeters2024wdc,li2020ditto}. This is a fair way to compare matchers, but it is not
what a catalogue experiences. Accepted pairs are merged transitively, so a single false link between two
product families moves every offer of one family into the other, while pairwise F1 counts that mistake
only once.

Between the matcher's scores and the final catalogue, several decisions are taken that the matcher does
not make itself. Someone has to pick the score above which offers are merged. Many pipelines also add
hand-written rules that block a merge when the offers disagree on brand, variant, capacity or model
number, because such rules are easy to audit. And when people can check only a fraction of the
decisions, someone has to choose which pairs they see. We ask how each of these choices behaves as the
underlying matcher gets better.

We study them on WDC Products~\cite{peeters2024wdc}, whose nine variants differ in the share of
difficult corner-case pairs and in the share of test products not seen during training, and on
\R{nextra} benchmarks from the Magellan collection~\cite{konda2016magellan,mudgal2018deepmatcher}. We
deliberately cover very different architectures, because a finding about thresholds or rules is only
useful if it does not depend on one model. The eight matchers include string similarity, the benchmark's
word co-occurrence baseline and gradient-boosted trees; a fine-tuned bi-encoder that embeds each offer
separately~\cite{reimers2019sbert,song2020mpnet}; two fine-tuned cross-encoders, RoBERTa~\cite{liu2019roberta}
and DeBERTa-v3~\cite{he2023debertav3}; and two zero-shot instruction-tuned LLMs from different
families~\cite{qwen2025,abdin2024phi3}. Nothing is tuned on test data. Before running any experiment we
wrote down how consistent across variants, and how significant, an effect must be before we report it as
a finding.

Our contributions are as follows.
\begin{itemize}
\item \IfClaim{B}{Choosing the threshold for entity-level F1 on validation, instead of for pairwise F1,
  improves the resolved catalogue of the cross-encoder in \R{CE.entbestwins} variants (entity F1
  \R{CE.entpw} to \R{CE.entbest}) while leaving the matcher unchanged.}{We compare pairwise and
  entity-level operating points under a validation-only protocol. Selecting for entity F1 on validation
  does not carry over to test reliably (\R{CE.entbestwins} variants for the cross-encoder), although a
  better test threshold exists (oracle \R{CE.entoracle}).}
\item \IfClaim{D}{A hard attribute veto loses precision as the matcher improves ($\rho=\R{rho}$,
  \R{npts} runs). A rule that removes errors from a weak matcher mostly removes correct merges from a
  strong one.}{We measure the precision of a hard attribute veto at each matcher's own threshold in
  \R{npts} runs and see no consistent link to matcher quality ($\rho=\R{rho}$).}
  \IfClaim{C}{With the cross-encoder it lowers AP by \R{CE.dapveto} (improved in \R{CE.dapvetowins}
  variants, $p=\R{CE.dapvetop}$).}{}
\item \IfClaim{F}{Conflicting attributes are more useful for directing review than for rejecting
  pairs: ranking pairs by a conflict-aware error probability recovers \R{CE.awareunc} more F1 than
  ranking by uncertainty at a \R{budget} budget.}{We test conflicting attributes as a signal for directing
  review and find no gain over calibrated uncertainty (\R{CE.awareunc} F1 at a \R{budget} budget), so
  uncertainty alone is enough to rank pairs for review.}
\item All code is public. One notebook rebuilds every table and figure from the public benchmarks, and
  our cross-encoder differs from the published RoBERTa numbers by \R{CE.paritybias} F1 on average over
  the nine variants.
\end{itemize}

% ======================================================================================================
\section{Related Work}

\textbf{Matchers and benchmarks.} Since Ditto~\cite{li2020ditto} showed that a pre-trained language
model fine-tuned on serialised record pairs beats earlier deep matchers, transformer cross-encoders have
been the usual strong baseline. Later work added an auxiliary entity-identification
objective~\cite{peeters2021jointbert} and contrastive pre-training on product
offers~\cite{peeters2022supcon}. Large language models can now match products with few or no labelled
pairs~\cite{peeters2023llm}; fine-tuning them improves both accuracy and transfer to unseen
products~\cite{steiner2025finetune}, and small models come close at a fraction of the
cost~\cite{zhang2024anymatch}. A controlled comparison still places cross-encoders ahead of
bi-encoders~\cite{zhang2026beyond}. WDC Products~\cite{peeters2024wdc} is our main benchmark and the
Magellan datasets~\cite{konda2016magellan,mudgal2018deepmatcher,kopcke2010evaluation} serve as a second
family. We do not propose a matcher; the matcher is the variable we change.

\textbf{After matching.} Errors become more expensive once pairs are merged into groups. GraLMatch uses
the structure of the match graph to find and remove wrong edges~\cite{dmeer2025gralmatch}, and TransClean
does so without labels by testing whether matches are transitively consistent~\cite{dmeer2025transclean}.
One of our candidate configurations is a gated variant of the TransClean idea, which we do not claim as
new. Other recent work clusters records directly with LLMs~\cite{fu2026llmcer} or spreads a small label
budget through the graph~\cite{wang2026alper}. Blocking limits recall before any of these steps
run~\cite{zeakis2023embeddings,moslemi2024blocking}; we use the candidate pairs released with each
benchmark so that our numbers stay comparable with earlier results.

\textbf{Confidence and review.} Neither a similarity score nor a classifier output is automatically a
probability. Platt scaling~\cite{platt1999probabilistic} and temperature scaling~\cite{guo2017calibration}
are the standard post-hoc corrections. For entity matching, LLM-based matchers have been found to be
somewhat overconfident~\cite{kamsteeg2025calibration}, and matchers built on decision models are well
calibrated only when the underlying model is strong~\cite{steiner2026decision}. We use calibrated
probabilities to choose pairs for review and ask whether attribute conflicts tell us anything that the
probabilities do not.

% ======================================================================================================
\section{Method}
\label{sec:method}

\subsection{Setting}
Each candidate pair $(r_i,r_j)$ receives a score $s_{ij}\in[0,1]$ from the matcher and has a gold label
$y_{ij}\in\{0,1\}$. We merge all pairs with $s_{ij}\ge\tau$ and take the connected components as the
resolved entities. To score them we use \emph{entity F1}: the pairwise F1 between the pairs implied by
the predicted clusters and those implied by the gold clusters, computed over the records that occur in
the evaluated split. A wrong link that joins two clusters is therefore charged for every pair it
creates.

\subsection{Calibration and operating points}
Raw scores are turned into probabilities with Platt scaling fitted on validation,
$p_{ij}=\sigma(a\,\mathrm{logit}(s_{ij})+b)$. We compare four ways of setting the operating point, all
of which look only at validation data:
(i)~$\tau_{\text{pw}}$ maximises pairwise F1, as most papers do;
(ii)~$\tau_{\text{ent}}$ maximises entity F1;
(iii)~$\tau_{\text{pw}}$ followed by gated transitive cleanup~\cite{dmeer2025transclean}: whenever a
record is linked to two others that are not linked to each other, the weaker of its two links is removed
if its calibrated probability is below a gate $\gamma\in\{0.4,0.6,0.8\}$, with $\gamma$ picked on
validation;
(iv)~the combination of $\tau$ and $\gamma$ that does best on validation.
We also show an oracle that picks the threshold on test, as an upper bound only.

\subsection{Attribute conflicts}
We reuse the constraint layer of our toolkit. For two offers, the conflict set $C_{ij}$ contains the
attributes that identify a product, namely brand, model line, variant, storage capacity and pack
quantity, whenever the two offers give different values for them. It also contains variant markers that appear in only one title, except when that title's tokens
include all tokens of the other one, which we treat as a shortened title. Finally, it contains numbers
with the same unit whose values disagree, such as \texttt{50mm} against \texttt{85mm}. A stricter version
keeps only conflicts that both offers state explicitly.

\emph{Hard veto.} Under the veto a pair is merged only if $s_{ij}\ge\tau_{\text{pw}}$ and
$C_{ij}=\varnothing$. Besides F1 at this shared threshold, we report the veto's precision. It is computed
over the pairs that the matcher accepts at $\tau_{\text{pw}}$ and the veto rejects, and gives the fraction
of them that are in fact non-matches. Pairs the matcher already rejects are left out, because blocking
them changes nothing. To see the effect on ranking, we cap the scores of conflicting pairs at $\kappa$,
$s'_{ij}=\min(s_{ij},\kappa)$, choosing $\kappa$ by average precision (AP) on validation.

\emph{Soft feature.} Rather than applying a fixed rule, a second calibrator receives the conflict
indicators together with the score,
$q_{ij}=\sigma(\mathbf{w}^\top[z_{ij},c^{f}_{ij},c^{s}_{ij},c^{f}_{ij}z_{ij},c^{s}_{ij}z_{ij}]+b)$, where
$z_{ij}=\mathrm{logit}(s_{ij})$ and $c^{f}$, $c^{s}$ mark conflicts under the full and the stricter
definition. It is fitted on validation and can therefore learn that some conflicts matter and others do
not.

\subsection{Review routing}
Suppose a reviewer can check $k=\beta N$ of the $N$ test pairs, with $\beta\in\{\R{budgets}\}$, and fixes
every decision they look at. Taking $d_{ij}=[s_{ij}\ge\tau_{\text{pw}}]$ as the automatic decision, the
chance that it is wrong is $e_{ij}=d_{ij}(1-p_{ij})+(1-d_{ij})\,p_{ij}$. \emph{Uncertainty routing}
gives the reviewer the $k$ pairs with the highest $e_{ij}$, and \emph{conflict-aware routing} does the
same with $q_{ij}$ in place of $p_{ij}$. For comparison we also place accepted pairs that have a conflict
at the front of the queue (\emph{conflict-first}) and select pairs at random.

\subsection{Statistics}
Every WDC variant is run with \R{nseeds} seeds. Each seed refits the matcher on its own \R{trainsub}
sample of the training pairs, so the spread we report reflects the training data and not only random
initialisation; results are averaged over seeds within a variant. We treat an effect as a finding when
it has the expected sign in at least \R{minwins} WDC variants and a two-sided Wilcoxon signed-rank test
over variants~\cite{wilcoxon1945} gives $p<\R{alpha}$. Individual runs come with 95\% intervals from a
paired Poisson bootstrap over test pairs (\R{nboot} resamples). These thresholds were set in the
experiment configuration before any result existed, and claims that miss them are reported as negative
findings.

% ======================================================================================================
\section{Experimental Setup}

\textbf{Data.} Table~\ref{tab:table_datasets} summarises the benchmarks. For WDC
Products~\cite{peeters2024wdc} we use the released pairs of the medium development set at 20, 50 and
80\% corner cases, each evaluated on test sets with 0, 50 and 100\% unseen products. The benchmark also
supplies gold product clusters and marks hard negatives. Abt-Buy, Amazon-Google~\cite{kopcke2010evaluation}
and Walmart-Amazon~\cite{konda2016magellan} are used with the splits common in the DeepMatcher line of
work~\cite{mudgal2018deepmatcher}. They have no gold clusters, so they enter only the pairwise analyses.

\GenTable{table_datasets}{Dataset statistics}

\textbf{Matchers.} \emph{Token-set} scores a pair by normalised token-set string similarity.
\emph{Word co-occurrence} marks the words that occur in both offers and trains a linear SVM with balanced
class weights; it is the simple baseline of the WDC Products paper~\cite{peeters2024wdc}. \emph{GBDT} is
histogram gradient boosting over string, token, numeric and attribute-agreement features. The
\emph{bi-encoder} starts from the MPNet sentence encoder all-mpnet-base-v2~\cite{song2020mpnet,reimers2019sbert}
and is fine-tuned on the same training pairs: each offer is mean-pooled on its own, and the match
probability is a learned logistic function of the cosine similarity. The two \emph{cross-encoders},
\R{cemodel}~\cite{liu2019roberta} and DeBERTa-v3-base~\cite{he2023debertav3}, read both offers together
and are fine-tuned for \R{ceepochs} epochs with batch size 32, learning rate $2\times10^{-5}$ and inputs
of up to 256 tokens. All three neural matchers serialise records field by field as in
Ditto~\cite{li2020ditto} and use early stopping on a slice held out from the training split. RoBERTa is
our primary strong matcher because published results exist for it on WDC Products; DeBERTa serves as an
independent replication. The two \emph{LLMs}, \R{llmmodel}~\cite{qwen2025} and
Phi-3.5-mini-instruct~\cite{abdin2024phi3}, are used without fine-tuning. Their score is
$P(\text{Yes})/(P(\text{Yes})+P(\text{No}))$ from a single forward pass over the prompt ``Do the following
two product offers refer to the same real-world product? Answer with Yes or No only.'', followed by both
titles and brands.

\textbf{Hardware.} All GPU work ran on one \R{gpu}. Training and scoring one cross-encoder took
\R{CE.minutes} minutes on average, and the LLM needed \R{LLM.mspair} ms per pair.

% ======================================================================================================
\section{Results}

\subsection{Matching quality and agreement with published results}

\GenTable{table_matching}{Matching quality}

The eight matchers in Table~\ref{tab:table_matching} range from AP \R{TS.ap} for token-set similarity to
\R{CE.ap} for the RoBERTa cross-encoder, with the bi-encoder at \R{BI.ap}, DeBERTa at \R{DB.ap} and the two
LLMs at \R{LLM.ap} and \R{LLMb.ap}. This spread is what allows us to relate the behaviour of thresholds
and rules to matcher quality. On the variant with 80\% corner cases and only seen products, the
cross-encoder reaches F1 \R{CE.fonehalfeighty} at its own decision boundary, against
\R{CE.pubeighty} published for RoBERTa-base~\cite{peeters2024wdc}. Over all nine variants our numbers
differ from the published ones by \R{CE.paritybias} on average (at most \R{CE.paritymax}), and those of
our word co-occurrence baseline by \R{WC.paritybias}. \IfClaim{P}{This is close enough to compare our
results directly with the published tables.}{This gap is larger than the tolerance we set, so any
comparison with published rows has to allow for it.} Neither LLM sees a single labelled pair. Calibration matters for the later sections: uncalibrated token-set scores
have an ECE of \R{TS.eceraw}, and after Platt scaling on validation the ECE is \R{TS.ececal} for
token-set and \R{CE.ececal} for the cross-encoder.

\subsection{Choosing the operating point for entities}
\label{sec:entity}

\GenTable{table_entity}{Entity-level F1}

Table~\ref{tab:table_entity} gives entity F1 under the four ways of setting the operating point.
\IfClaim{B}{Choosing the threshold and the cleanup gate for entity F1 on validation improves the
cross-encoder from \R{CE.entpw} to \R{CE.entbest}, with a gain in \R{CE.entbestwins} variants
($p=\R{CE.entbestp}$). For GBDT, entity F1 rises from \R{GB.entpw} to \R{GB.entbest}
(\R{GB.entbestwins} variants), and its largest cluster shrinks on average from \R{GB.maxclpw} to
\R{GB.maxclbest} offers; most of the gain therefore comes from fewer chains of wrong links between
unrelated products. The method needs no labels beyond the validation split that the protocol uses
anyway.}{Tuning the threshold for entity F1 on validation does not reliably help on test. The
cross-encoder gains in only \R{CE.entbestwins} variants ($p=\R{CE.entbestp}$; entity F1 \R{CE.entpw}
to \R{CE.entbest}). The test-tuned oracle (\R{CE.entoracle}) shows that a better threshold exists, but
the validation split is too small to locate it with confidence. We read this as a warning against tuning
entity-level settings on small development sets.} The oracle column indicates how much room each
matcher leaves; for GBDT it reaches \R{GB.entoracle}, against \R{GB.entpw} at the pairwise threshold.

\subsection{Hard constraints and matcher quality}
\label{sec:veto}

\GenTable{table_constraint}{Hard veto vs.\ soft conflict feature}
\GenFig{figure}{fig_veto_vs_quality}{\columnwidth}{Precision of the hard veto at each matcher's own
threshold, plotted against the matcher's test AP. Each point is one matcher, variant and seed; marker
shapes distinguish the benchmarks. Below 0.5 the veto rejects more true matches than
non-matches.}{fig:veto}

\IfClaim{D}{Fig.~\ref{fig:veto} shows the main pattern: the stronger the matcher, the less precise the
veto (Spearman $\rho=\R{rho}$, $p=\R{rhop}$, $n=\R{npts}$). Of the accepted pairs that the veto would
reject, \R{TS.vetoprec} are non-matches for token-set similarity, \R{GB.vetoprec} for GBDT,
\R{BI.vetoprec} for the bi-encoder, and only \R{CE.vetoprec} and \R{DB.vetoprec} for the RoBERTa and
DeBERTa cross-encoders. A strong matcher has already turned down most pairs whose
attributes really differ. What is left for a rule to overrule are mostly true matches whose offers
describe the product in different words, for example through a shortened title or a seller's own naming.
The Magellan benchmarks give the same picture (cross-encoder veto precision
\R{CE.genvetoprec}).}{Fig.~\ref{fig:veto} shows no consistent link between how strong a matcher is and
how precise the veto is ($\rho=\R{rho}$, $p=\R{rhop}$). Veto precision is \R{TS.vetoprec},
\R{GB.vetoprec} and \R{CE.vetoprec} for token-set, GBDT and the cross-encoder.}
\IfClaim{C}{On the cross-encoder the veto reduces AP by \R{CE.dapveto} ($p=\R{CE.dapvetop}$) and F1 at
the shared threshold by \R{CE.dfoneshared}, blocking \R{CE.trueblocked} true matches per variant.}{Its
effect on the cross-encoder's AP (\R{CE.dapveto}; improved in \R{CE.dapvetowins} variants,
$p=\R{CE.dapvetop}$) is not consistent across the grid.}
If the veto only uses conflicts that both offers state, its precision on the cross-encoder becomes
\R{CE.vetoprecsym}, and AP changes by \R{CE.dapsym}.

\IfClaim{E}{Given to a calibrator as features, the same conflicts help rather than hurt: AP changes by
\R{GB.dapsoft} for GBDT (\R{GB.dapsoftwins} variants) and by \R{CE.dapsoft} for the cross-encoder
(\R{CE.dapsoftwins}). Because its weights are learned, the calibrator can discount conflicts that turn
out not to matter, which a fixed rule cannot do.}{Given to a calibrator as features, the conflicts change
AP by \R{GB.dapsoft} for GBDT (\R{GB.dapsoftwins} variants) and by \R{CE.dapsoft} for the cross-encoder
(\R{CE.dapsoftwins}). They add little ranking information beyond what the matchers already capture.}

\subsection{Conflict-aware review routing}
\label{sec:review}

\GenFig{figure*}{fig_review}{0.85\textwidth}{F1 gained over no review as a function of the share of test
pairs a reviewer corrects (WDC Products; mean and standard deviation over variants).}{fig:review}
\GenTable{table_review}{Review routing}

Once probabilities are calibrated, spending a \R{budget} review budget on the most uncertain decisions
adds \R{CE.revunc} F1 for the cross-encoder and \R{GB.revunc} for GBDT, whereas random selection adds
\R{CE.revrandom} and \R{GB.revrandom} (Fig.~\ref{fig:review}).
\IfClaim{F}{Including the conflicts in the routing model helps further, by \R{CE.awareunc} F1 for the
cross-encoder (\R{CE.awareuncwins} variants, $p=\R{CE.awareuncp}$) and \R{GB.awareunc} for GBDT. The
conflict signal thus changes role. As an automatic veto it hurts a strong matcher, but as a reason to
show a pair to a person it is useful, since a reviewer can tell a real conflict from a naming difference
and a rule cannot. The effect also appears on the Magellan benchmarks for the cross-encoder
(\R{CE.genawareunc}, positive in \R{CE.genawarewins} datasets).}{Including the conflicts in the routing
model does not improve on uncertainty routing (cross-encoder \R{CE.awareunc}, \R{CE.awareuncwins}
variants, $p=\R{CE.awareuncp}$). After calibration, the conflicts carry no extra information about which
decisions are wrong.} \IfClaim{R}{The DeBERTa cross-encoder leads to the same conclusions as RoBERTa
for both the veto (AP \R{DB.dapveto}) and review routing (\R{DB.awareunc} F1).}{The DeBERTa
cross-encoder does not fully agree with RoBERTa (veto AP \R{DB.dapveto}, review routing \R{DB.awareunc}
F1), so these two findings should be read as specific to the primary matcher.}
Table~\ref{tab:table_review} also reports conflict-first routing, which is the most direct way of
turning the veto into a review trigger.

\subsection{Generalisation and robustness}

\GenTable{table_generalisation}{Generalisation to the Magellan benchmarks}
\GenTable{table_robustness}{Robustness to unseen products}

Table~\ref{tab:table_generalisation} repeats the pairwise analyses on the Magellan benchmarks. There the
cross-encoder reaches AP \R{CE.genap}; the veto changes its AP by \R{CE.gendapveto} and the soft
calibrator by \R{CE.gendapsoft}. Table~\ref{tab:table_robustness} shows how F1 changes when the test
products were not seen in training. Going from seen to unseen products, the cross-encoder loses
\R{CE.unseendrop} F1, the LLM \R{LLM.unseendrop} and word co-occurrence \R{WC.unseendrop}. WDC
Products was designed to expose this kind of drop~\cite{peeters2024wdc}.

% ======================================================================================================
\section{Discussion}

\textbf{Practical advice.} Three recommendations follow from our results. First,
\IfClaim{B}{set the threshold for the output that is actually used; tuning it for entity F1 on
validation costs nothing extra.}{if entity quality matters, check the entity-level metric on validation,
but expect the choice to be noisy when the development set is small.} Second, \IfClaim{D}{be careful with
hard attribute rules on top of a strong matcher, because their precision drops as the matcher
improves.}{measure a hard rule's precision at the threshold that will actually be deployed before
adopting it.} Third, choose pairs for review by their calibrated error
probability\IfClaim{F}{, and include attribute conflicts in that model.}{.}

\textbf{Threats to validity.} Our reviewer is simulated and always returns the gold answer; real
reviewers make mistakes, and some pairs take longer to check than others. Both LLMs are used zero-shot
with a single prompt. Better prompts or fine-tuning~\cite{steiner2025finetune} would make them stronger
and, if the pattern in Section~\ref{sec:veto} holds, would also make the veto less precise for them. Entity-level
results come from WDC Products alone, because the Magellan benchmarks contain no gold clusters. The
attribute extractor relies on rules and a general lexicon rather than one tailored to each product
category. Finally, we evaluate on the candidate pairs released with each benchmark, which places blocking
outside the scope of this study.

% ======================================================================================================
\section{Conclusion}

We studied what happens between a matcher's pair scores and the resolved catalogue, using eight matchers
from four architecture families on WDC Products and \R{nextra} Magellan benchmarks. \IfClaim{B}{Tuning the operating point for entities
on validation improved the merged result.}{The pairwise operating point proved hard to beat with
validation data alone.} \IfClaim{D}{A hard attribute veto lost precision as the matcher it was meant to
protect became better,}{The value of a hard attribute veto did not follow matcher quality,} and
\IfClaim{F}{the same conflict evidence was most useful for deciding which pairs a person should
check.}{calibrated uncertainty was enough to decide which pairs a person should check.} Every number in
this paper can be regenerated from public data with a single notebook.

\section*{Data and Code Availability}
Code, the experiment notebook and all per-pair scores: \url{https://github.com/Vikramsrini/MajorProject}.
WDC Products: \url{https://webdatacommons.org/largescaleproductcorpus/wdc-products/}.

\section*{Acknowledgment}
% IEEE policy: the use of AI-generated content (text, code, figures, images) must be disclosed here.
% Edit the second sentence so that it states accurately what was AI-assisted; do not delete it.
[Thank your guide, department and any compute support here.]
{\color{red}[Required AI-use disclosure, edit before submission:]} Parts of the experiment code and early
drafts of this manuscript were prepared with the help of an AI assistant (Claude, Anthropic). The authors
reviewed, tested and revised all code and text and take full responsibility for the content.

\bibliographystyle{IEEEtran}
\bibliography{refs}

\end{document}
